\color{unimain}{\subsection{Контроль силового выключателя}}
\color{black}

\begin{enumerate}

%%%%%%%%%%%%%%%%%%%%%%%%%%%%%%%%%%%%%%%%%%%%%%%%%%%%%%%%%%%%%%%%%%%% Контроль В1 %%%%%%%%%%%%%%%%%%%%%%%%%%%%%%%%%%%%%%%%%%%%%%%%%%%%%%%%%%%%%%%%%%%%

\item
Функциональный блок контроля силового выключателя (КСВ) предназначен для обработки
технологической информации от выключателя и его цепей управления, а также от трансформатора тока, позволяя определять их состояние,
информировать о возможных неисправностях и формировать соответствующую сигнализацию.

\item
В случае подключения присоединения к распределительному устройству через два выключателя для
каждого из них можно задействовать свою логику КСВ.

\item
Параметры для настройки КСВ представлены в таблице \ref{app_set:ksv}.

\item
Обобщенная функциональная схема логики контроля выключателя приведена на рисунке \ref{pict_ksv:KSV_control}.

\medskip
\begin{figure}[H]
\centering
\includegraphics[width=0.8\textwidth,height=0.8\textheight,keepaspectratio]{\fbpath/10.02 АУВ/КСВ/img/KSV_control.pdf}
\captionof{figure}{Функциональная схема логики контроля выключателя}\label{pict_ksv:KSV_control}
\end{figure}

\item
Функция КСВ осуществляет контроль исправности цепей управления выключателем. Контроль
целостности цепей управления осуществляется по сигналам состояния цепей первой и второй группы
электромагнитов отключения (ЭМО1 и ЭМО2) и цепи электромагнита включения (ЭМВ). Количество групп
электромагнитов отключения определяется положением программного переключателя \textbf{«Кол-во ЭМО»}. При
одновременном наличии или отсутствии сигналов контроля цепей ЭМО и ЭМВ в течение времени,
превышающем 10 с, формируется сигнал о неисправности цепей управления. 

\item
Обнаружение неисправности выключателя осуществляется по следующим признакам:
\begin{itemize}
	\item низкое или аварийное давление элегаза (SF6) в выключателе (с задержкой 1 с);
	\item отключенное положения автомата ШП привода выключателя (с задержкой 0,1 с);
	\item неисправность обогрева выключателя (с задержкой 1 с);
	\item неисправность привода – при отсутствии заведенного состояния пружины (с задержкой 20 с);
	\item неисправность цепей управления;
	\item срабатывание ЗНФ;
	\item неисправность цепей ОТ ЭМО1 или ЭМО2 (при его наличии).
\end{itemize}


%%%%%%%%%%%%%%%%%%%%%%%%%%%%%%%%%%%%%%%%%%%%%%%%%%%%%%%%%%%%%%%%%%%% Контроль ТТ В1 %%%%%%%%%%%%%%%%%%%%%%%%%%%%%%%%%%%%%%%%%%%%%%%%%%%%%%%%%%%%%%%%%%%%

\item
Предусмотрен контроль состояния ТТ. Обеспечивается прием сигналов о низком и аварийном снижении
давлении элегаза в ТТ. Предусматривается действие на отключение выключателя при аварийном снижении
давления элегаза в ТТ. Отключение при аварийном снижении элегаза в ТТ обеспечивается при включенном
положении программной накладки \textbf{«Откл.авар.давл.SF6 ТТ»}.

\item
Функциональная схема логики контроля ТТ выключателя приведена на рисунке \ref{pict_ksv:KSV_control_TT_B1}.

\medskip
\begin{figure}[H]
\centering
\includegraphics[width=0.8\textwidth,height=0.8\textheight,keepaspectratio]{\fbpath/10.02 АУВ/КСВ/img/KSV_control_CT.pdf}
\captionof{figure}{Функциональная схема логики контроля ТТ выключателя}\label{pict_ksv:KSV_control_TT_B1}
\end{figure}


%%%%%%%%%%%%%%%%%%%%%%%%%%%%%%%%%%%%%%%%%%%%%%%%%%%%%%%%%%%%%%%%%%%% Защита ЭМУ В1 %%%%%%%%%%%%%%%%%%%%%%%%%%%%%%%%%%%%%%%%%%%%%%%%%%%%%%%%%%%%%%%%%%%%

\item
Предусмотрена защита электромагнитов выключателя от повреждений при длительности протекании
тока. Контроль длительности протекания токов через электромагниты управления осуществляется с
использованием датчиков тока, выходные контакты которых подключены в соответствующим входам
устройства \textit{«Работа ЭМВ»}, \textit{«Работа ЭМО1»}, \textit{«Работа ЭМО2»}. При протекании через электромагнит
управления тока в течение времени, превышающем регулируемую задержку \textbf{«Тср защ.ЭМУ»} формируется
выходной сигнал, действующий на соответствующий независимый расцепитель защитного автомата питания
цепей электромагнитов.

\item
Обобщенная функциональная схема логики защиты ЭМУ выключателя приведена на рисунке \ref{pict_ksv:KSV_EMU_protection}.

\medskip
\begin{figure}[H]
\centering
\includegraphics[width=0.8\textwidth,height=0.8\textheight,keepaspectratio]{\fbpath/10.02 АУВ/КСВ/img/KSV_EMU_protection.pdf}
\captionof{figure}{Функциональная схема логики защиты ЭМУ выключателя}\label{pict_ksv:KSV_EMU_protection}
\end{figure}


%%%%%%%%%%%%%%%%%%%%%%%%%%%%%%%%%%%%%%%%%%%%%%%%%%%%%%%%%%%%%%%%%%%% Блокировка управления В1 %%%%%%%%%%%%%%%%%%%%%%%%%%%%%%%%%%%%%%%%%%%%%%%%%%%%%%%%%%%%%%%%%%%%

\item 
Предусмотрена блокировка отключения выключателя при аварийном снижении элегаза или наличии внешнего сигнала блокировки управления.

\item
Включение выключателя блокируется в следующих случаях:
\begin{itemize}
	\item ремонт выключателя;
	\item отключенное положения автомата ШП привода выключателя;
	\item неисправность цепей управления выключателем;
	\item аварийное снижение давления элегаза в выключателе;
	\item внешняя блокировка включения или управления выключателем;
	\item отсутствие разрешения на включение от САУ (при включенной накладке «Разр.вкл.САУ»);
	\item ожидание разряда БСК (при включенной накладке \textbf{«Конт.разр.БСК»} после отключения БСК запрещается ее включение в течение времени \textbf{«Тожид разряда БСК»}).
\end{itemize}

\item
Обобщенная функциональная схема логики блокировки управления выключателя приведена на рисунке \ref{pict_ksv:KSV_control_block}.

\medskip
\begin{figure}[H]
\centering
\includegraphics[width=0.8\textwidth,height=0.8\textheight,keepaspectratio]{\fbpath/10.02 АУВ/КСВ/img/KSV_control_block.pdf}
\captionof{figure}{Функциональная схема логики блокировки управления выключателя}\label{pict_ksv:KSV_control_block}
\end{figure}

%%%%%%%%%%%%%%%%%%%%%%%%%%%%%%%%%%%%%%%%%%%%%%%%%%%%%%%%%%%%%%%%%%%% Реле фиксации В1 %%%%%%%%%%%%%%%%%%%%%%%%%%%%%%%%%%%%%%%%%%%%%%%%%%%%%%%%%%%%%%%%%%%%

\item
В устройстве предусмотрена программная фиксация команд или включенного положения выключателя
в зависимости от положения программного переключателя \textbf{«Реле фиксации»}. В положении \textit{«РФК»}
фиксируется команда включения, в положении \textit{«РФП»} – включенное положение. Квитирование сигнала
фиксаци может осуществляться от сброса сигнализации или подачей оперативной команды на отключение.
При отключении выключателя от защит и наличии сигнала \textit{«Фиксация»} формируется сигнал о факте аварийного
отключения выключателя \textit{«Сигнал несоотв.»}, необходимый для работы АПВ.

\item
Обобщенная функциональная схема логики реле фиксации выключателя приведена на рисунке \ref{pict_ksv:KSV_relay_fix}.

\medskip
\begin{figure}[H]
\centering
\includegraphics[width=0.8\textwidth,height=0.8\textheight,keepaspectratio]{\fbpath/10.02 АУВ/КСВ/img/KSV_relay_fix.pdf}
\captionof{figure}{Функциональная схема логики реле фиксации выключателя}\label{pict_ksv:KSV_relay_fix}
\end{figure}

\end{enumerate}